W~baggingu i~lesie losowym każde drzewo decyzyjne uczy się na osobnej losowej próbce danych treningowych, a~jakość zespołu zależy od dokładności drzew i~od tego, jak bardzo różnią się ich predykcje. Obie wielkości mierzy się zwykle dopiero po uczeniu. W~pracy sprawdziliśmy, czy proste miary różnorodności i~reprezentatywności samych próbek, obliczone przed uczeniem, są związane z~jakością zespołu. Porównaliśmy 15 sposobów losowania próbek: losowanie ze zwracaniem i~bez zwracania o~różnym rozmiarze, losowanie próbek o~wspólnej części przy stałym rozmiarze oraz losowanie warstwowe. Badanie objęło 12 zbiorów klasyfikacyjnych i~12 regresyjnych, pięciokrotną walidację krzyżową powtórzoną 10 razy i~łącznie 29\,400 zespołów po 100 drzew. Jednostką w~testach statystycznych był zbiór danych. Żaden sposób losowania nie dał lepszych zespołów niż losowanie stosowane domyślnie, czyli losowanie ze zwracaniem próbki o~rozmiarze części treningowej. Istotnie gorsze były próbki w~dużej części wspólne dla wszystkich drzew: przy tym samym rozmiarze próbek losowanie niezależne dało lepszy zespół w~21 z~24 zbiorów. Jakość zespołu zależała od liczby różnych obserwacji w~próbce, a~nie od tego, czy losowano ze zwracaniem. Miary próbek były istotnie, ale bardzo słabo związane z~różnicami między drzewami (współczynnik korelacji rang około 0.02--0.03). Po uwzględnieniu dokładności drzew większa różnorodność próbek wiązała się z~lepszym zespołem. Żadna pojedyncza miara nie wskazała jednak najlepszego sposobu losowania ani lepszego z~kilku losowań tym samym sposobem. Wyjaśniamy to tym, że przy ustalonym sposobie losowania miary próbek prawie się nie zmieniają, a~ich rozrzut maleje wraz z~liczbą obserwacji.
